\documentclass[12pt]{article}

%\usepackage[utf8]{inputenc}
%\usepackage[T2A]{fontenc}
\usepackage[russian,english]{babel}
\usepackage{amsmath,amsfonts,amssymb,bm}

\setlength{\textwidth}{168mm}%210
\setlength{\textheight}{210mm}%297
\topmargin -10mm
\oddsidemargin  -1mm
\evensidemargin -1mm
\usepackage[utf8]{inputenc}
\usepackage{amsmath,amssymb}
\usepackage{latexsym}
\usepackage{slashed}
\usepackage{authblk}
\usepackage{xcolor}
\usepackage{hyperref}

\usepackage{cancel}
\usepackage{indentfirst}
\usepackage{graphicx}% Include figure filesF
\usepackage{epstopdf}
%\graphicspath{{figs/}}
\usepackage{dcolumn}% Align table columns on decimal point
\usepackage{bm}% bold math
\usepackage{lipsum}
\usepackage{slashed}
\usepackage{enumitem}
\usepackage{txfonts}
\usepackage{esint}
\usepackage{xcolor}
\usepackage{mathtools}
\usepackage{transparent}
\usepackage{upgreek}
\usepackage{tabularx}
\usepackage{csquotes}
\usepackage{siunitx}
%\usepackage{geometry}
%\geometry{a4paper,total={178mm,257mm},right=15mm,left=15mm,bottom=20mm,top=20mm}
%\usepackage[font=footnotesize,labelfont=bf,labelsep=period]{caption}
% \usepackage[bookmarks=true,colorlinks=true, allcolors=black,linkcolor=black]{hyperref}
%\usepackage{tikz}
%\usepackage[compat=1.1.0]{tikz-feynman}
%\tikzfeynmanset{/tikzfeynman/warn luatex=false}
%\usetikzlibrary{arrows,automata}
%\usetikzlibrary{matrix}
%\addto\captionsenglish{\renewcommand{\figurename}{Fig.}}

%\pagestyle{plain}%номера страниц
%\usepackage{setspace}%Полуторный
%\onehalfspacing% межстрочный интервал
%\doublespacing

%\usepackage{wrapfig}
% \usepackage[backend=biber,style=ieee,articletitle=false]{biblatex}
%%\usepackage[backend=biber, style=phys,articletitle=false,biblabel=brackets,
%%            chaptertitle=false,pageranges=false]{biblatex}
%%\ExecuteBibliographyOptions[misc]{eprint=true}
%%\ExecuteBibliographyOptions[online]{eprint=true}
%%\addbibresource{main.bib}

%\usepackage{definitions}

\begin{document}
	\title{Hawking Radiation and Primordial Black Holes}
	
	\author{Borodin Nikita}
	
	
	
	\maketitle
	\newpage
	\section{Introduction}
	
	\begin{itemize}
		\item  \(e^{i \vec{k} \cdot \vec{x}}\)
		\item  \(\frac{e^{i k \cdot r}}{r}\)\\
	\end{itemize}
	
	\section{Find wavefunction}
	\indent Using the Fourier transform
	\[
	\begin{aligned}
		& \psi\left(\vec{r}_{1}, \vec{r}_{2}\right)=\int \frac{d^{3} p_{1} d^{3} p_{2} \psi\left(\vec{p}_{1} \vec{p}_{2}\right)}{\left.(2 \pi)^{6} \sqrt{2 E_{p} r_{1}}\right) } e^{i(...)} \\
		& =\int \frac{d^{2} p_1 d^{2} p_2}{\left(2 \pi \right)^3 } \langle f \overline{f} \mid S-1\rangle e^{i(...)}\\
		& =\int \frac{d^{3}p}{(2 \pi^3)} \frac{1}{2 E_p} \frac{m_f}{v} \overline{u}(\vec{p_1},\vec{s_1}) u(-\vec{p_2},\vec{s_2}) e^{i(...)}
	\end{aligned}
	\]
	
	\[
	\begin{aligned}
		S \mid i \rangle = \mid f \rangle = ...+ \int \frac{d^3 p_1 d^3 p_2}{(2 \pi)^6} \frac{\psi\left( \vec{p_1}, \vec{p_2}\right) }{\sqrt{2 E_1 2 E_2}} \mid \vec{p_1}, \vec{s_1}; \vec{p_2}, \vec{s_2} \rangle + ...
	\end{aligned}
	\]
	
	
	\[
	\begin{aligned}
		\langle \vec{p}_1 \vec{p}_2 \mid \vec{p}_1 \vec{p}_2 \rangle = (2 \pi)^3 \sqrt{2 E_1 2 E_2} \left[ \delta^{3}\left(\vec{p}_{1}-\vec{p}_{2}\right) \delta^{\prime}\left(\vec{p}_{2}-\vec{p}_{1}^{\prime}\right) - \delta^{3}\left(\vec{p}_{1}-\vec{p}_{2}\right) \delta^{\prime}\left(\vec{p}_{2}-\vec{p}_{1}^{\prime}\right) \right]
	\end{aligned}
	\]
	
	Momentum 
	
	\[
	\begin{aligned}
		&\mid p \rangle \equiv \sqrt{2 E_p} a_p^{\dagger} \mid 0 \rangle \\
		\langle \vec{q} \mid \vec{p} \rangle &= \langle \vec{q} \mid a_q \sqrt{2 E_q} \sqrt{2 E_p} a^{\dagger}_p \mid 0 \rangle = 2 E_p \left( 2 \pi \right)^3 \delta^3(\vec{q} - \vec{p}) \\
		&a_q a^{\dagger}_p = -a_q^{\dagger} a_p + \underbrace{\left\lbrace a_q, a^{\dagger}_p\right\rbrace }_{\left( 2 \pi^3\right) \delta^3 \left( \vec{q} - \vec{p}\right)  }
	\end{aligned}
	\]
	
	We'll get the following norm
	\[
	\begin{aligned}
		\langle \vec{q}, s \mid \vec{p}, r \rangle = 2 E_p \left(2 \pi \right)^3 \delta^3 \left( \vec{q} - \vec{p}\right) \delta_{sr}  
	\end{aligned}
	\]
	
	Substitute in our decay
	\[
	\begin{aligned}
		\langle q_1, r_1; q_2, r_2 \mid p_1 s_1; p_2 s_2 \rangle &= \sqrt{2 E_{p_1} 2 E_{p_2} 2 E_{q_1} 2 E_{q_2}} \langle 0 \mid a_{q_1 r_1} a_{q_2 r_2} a_{p_1 s_1}^{\dagger} a_{p_2 s_2}^{\dagger} \mid 0 \rangle \\
		&=\sqrt{2 E_{p_1} 2 E_{p_2} 2 E_{q_1} 2 E_{q_2}} \langle 0 \mid a_{q_1 r_1} \left(-a_{q_2 r_2}^{\dagger} a_{p_1 s_1} + \left\lbrace a_{q_2 r_2}, a^{\dagger}_{p_1 s_1}\right\rbrace \right)  a_{p_2 s_2}^{\dagger} \mid 0 \rangle \\
		&= \left\lbrace a_{q_2 r_2}, a^{\dagger}_{p_1 s_1}\right\rbrace \langle 0 \mid a_{q_1 r_1} a_{p_2 s_2}^{\dagger} \mid 0 \rangle - \langle 0 \mid a_{q_1 r_1} a_{q_2 r_2}^{\dagger} a_{p_1 s_1} a_{p_2 s_2}^{\dagger} \mid 0 \rangle \\
		&=\left\lbrace a_{q_2 r_2}, a^{\dagger}_{p_1 s_1}\right\rbrace \left\lbrace a_{q_1 r_1}, a^{\dagger}_{p_2 s_2}\right\rbrace - \left\lbrace a_{q_1 r_1}, a^{\dagger}_{q_2 r_2}\right\rbrace \left\lbrace a_{p_1 s_1}, a^{\dagger}_{p_2 s_2}\right\rbrace\\
	\end{aligned}
	\]
	
	Finally will get
	\[
	\begin{aligned}
		\langle \vec{q}_1 r_1; \vec{q}_2 r_2 \mid S \mid H \rangle &= \int \frac{d^3 p_1 d^3 p_2}{(2 \pi)^6 \sqrt{2 E_{p_1} 2 E_{p_2}}} \widetilde{\psi}(\vec{p_1}, \vec{p_2}) \langle q_1 r_1; q_2 r_2 \mid p_1 s_1; p_2 s_2\rangle \\
		&= \int \frac{d^3 p_1 d^3 p_2}{\left( 2 \pi\right) ^6 \sqrt{2 E_{p_1} 2 E_{p_2}}} \widetilde{\psi}(\vec{p_1}, \vec{p_2}) \left[ \left( 2 \pi\right) ^6 \delta^{3} (q_2 - p_1) \delta_{r_2 s_1} \delta^{3}(q_1 - p_1) \delta_{r_2 s_1} - \left( 2 \pi\right) ^6 \delta^{3} (q_2 - p_2) \delta_{r_2 s_2} \delta^{3} (q_1 - p_1) \delta_{r_1 s_1} \right] \\
		&= \frac{\widetilde{\psi}(q_2,q_1)}{\sqrt{2 E_{q_1} 2 E_{q_2}}} \delta_{r_2 s_1} \delta_{r_1 s_2} - \frac{\widetilde{\psi}(q_1,q_2)}{\sqrt{2 E_{q_1} 2 E_{q_2}}} \delta_{r_1 s_1} \delta_{r_2 s_2}
	\end{aligned}
	\]
	??
	\[
	\begin{aligned}
		\frac{\widetilde{\psi}(q_2 r_2;q_1 r_1) - \widetilde{\psi}(q_1 r_1;q_2 r_2)}{\sqrt{2 E_{q_1} 2 E_{q_2}}} 
	\end{aligned}
	\]
	
	\[
	\begin{aligned}
		\widetilde{\psi}(q_1 r_1; q_2 r_2) \sim \sqrt{2 E_{q_1} 2 E_{q_2}} \langle q_1 r_1; q_2 r_2 \mid S - 1 \mid H \rangle
	\end{aligned}
	\]
	
	\[
	\begin{aligned}
		\widetilde{\psi}(\vec{r_1} s_1; \vec{r_2} s_2) &= \int \frac{d^3 p_1 d^3 p_2}{(2 \pi)^6 \sqrt{2 E_{p1} 2 E_{p2}}} \widetilde{\psi} e^{i\vec{p_1} \cdot \vec{r_1} + i \vec{p_1} \vec{r_2}} \\
		&= \int \frac{d^3 p_1 d^3 p_2}{(2 \pi)^6} \underbrace{\langle p_1 s_1; p_2 s_2\mid S-1 \mid H\rangle}_{\left( 2 \pi \right)^4 \delta^4 \left(P_H - P_1 - P_2 \right) i \frac{m_f}{v} \overline{u}(p_1 s_1) v(p_2 s_2) } e^{i\vec{p_1} \cdot \vec{r_1} + i \vec{p_1} \vec{r_2}}
	\end{aligned}
	\]
	
	Using the spherical system of coordinates
	\[
	\begin{aligned}
		\int \frac{d^3 p_1 d^3 p_2}{(2 \pi)^6} \left( 2 \pi \right)^4 \delta^4 \left(P_H - P_1 - P_2 \right) i \frac{m_f}{v} \overline{u}(p_1 s_1) v(p_2 s_2) \Big|_{\vec{p_2} = -\vec{p_1}} e^{i\vec{p_1}  \cdot \vec{r_1} + i \vec{p_1} \vec{r_2}}
	\end{aligned}
	\]
	
	\[
	d^3 P_1 = \underbrace{P_1^2 dP_1}_{p_1 E_1 dE_1} d \Omega_{p_1} 
	\]
	
	\section{Fermion Antifermion product}
	From Dirac equation for fermion and antifermion 
	\section{Conclusion}
	\indent In this work, the various masses of the primordial black holes and their lifetimes were calculated. The work also investigated what would happen if there were a  black hole in the center of the Earth. It turned out that it would not be able to absorb all the matter around it. The maximum mass of a black hole that can have and still evaporate in the presence of the cosmic microwave background was calculated. The photon and neutrino energies of the primary black holes were calculated and compared with the range of measured energies of existing experimental facilities. It turned out that we need a new generation of installations with better resolution to study photons and neutrinos.
	
	
	%\LaTeX{} \cite{MacGibbon_1998}
	%\bibliographystyle{unsrt}
	%\bibliography{ref}
	
	\LaTeX{} \cite{MacGibbon_1998} is a set of macros built atop \TeX{} \cite{MacGibbon_1998}.
	\bibliographystyle{plain} % We choose the "plain" reference style
	\bibliography{ref} % Entries are in the refs.bib file
	
\end{document}